\documentclass[12pt]{article}
\usepackage{amsmath,amssymb,amsfonts,amsthm}
\usepackage[margin=1in]{geometry}

\begin{document}

\begin{center}
{\Large\bfseries Chapter 4, Problem 16}\\[6pt]
Byron \& Fuller, \textit{Mathematics of Classical and Quantum Physics}
\end{center}

\bigskip

\textbf{Problem.}
\begin{itemize}
\item[(a)] Find the eigenvalues and eigenvectors of the Lorentz transformation. Verify that the eigenvectors are null vectors.
\item[(b)] Evaluate the determinant of the pure Lorentz matrix.
\item[(c)] Any other matrix representation of the Lorentz transformation can be found by performing a similarity transformation on the matrix of Problem 14. What is the value of the determinant of any matrix representation of the Lorentz transformation?
\end{itemize}

\bigskip
\textbf{Solution.}

\textbf{(a)} Consider the standard Lorentz boost in the $x$-direction (working in the $ct$-$x$ subspace):
\[
L = \begin{pmatrix} \gamma & -\beta\gamma \\ -\beta\gamma & \gamma \end{pmatrix}.
\]

The eigenvalue equation $\det(L - \lambda I) = 0$ gives:
\[
(\gamma - \lambda)^2 - \beta^2\gamma^2 = 0 \implies \gamma - \lambda = \pm\beta\gamma.
\]
\[
\lambda_1 = \gamma(1-\beta) = \gamma\cdot\frac{1-\beta}{1} = \sqrt{\frac{1-\beta}{1+\beta}}, \quad \lambda_2 = \gamma(1+\beta) = \sqrt{\frac{1+\beta}{1-\beta}}.
\]

Note: $\lambda_1\lambda_2 = 1$, consistent with $\det L = 1$.

\textbf{Eigenvectors:} For $\lambda_1 = \gamma(1-\beta)$:
\[
(\gamma - \lambda_1)\,v_1 - \beta\gamma\,v_2 = 0 \implies \beta\gamma\,v_1 = \beta\gamma\,v_2 \implies v_1 = v_2.
\]
So $\mathbf{e}_1 = \begin{pmatrix} 1 \\ 1 \end{pmatrix}$, corresponding to $(ct, x)$ with $ct = x$.

For $\lambda_2 = \gamma(1+\beta)$:
\[
(\gamma - \lambda_2)\,v_1 - \beta\gamma\,v_2 = 0 \implies -\beta\gamma\,v_1 = \beta\gamma\,v_2 \implies v_1 = -v_2.
\]
So $\mathbf{e}_2 = \begin{pmatrix} 1 \\ -1 \end{pmatrix}$, corresponding to $ct = -x$.

\textbf{Verification that eigenvectors are null:} Using the Minkowski metric $ds^2 = c^2t^2 - x^2$:
\begin{align*}
\text{For } \mathbf{e}_1: &\quad (ct)^2 - x^2 = 1 - 1 = 0. \checkmark \\
\text{For } \mathbf{e}_2: &\quad (ct)^2 - x^2 = 1 - 1 = 0. \checkmark
\end{align*}
Both eigenvectors lie along the light cone and are therefore null vectors.

\bigskip
\textbf{(b)} As computed in Problem 12(a):
\[
\det L = \gamma^2(1-\beta^2) = \gamma^2 \cdot \frac{1}{\gamma^2} = 1.
\]

\bigskip
\textbf{(c)} The determinant is invariant under similarity transformations:
\[
\det(P^{-1}LP) = \det(P^{-1})\det(L)\det(P) = \det(L) = 1.
\]
Therefore the determinant of any matrix representation of a proper Lorentz transformation is $+1$. \qed

\end{document}
